% =====================================================================
% INTRODUCTION — NMI prints this with no heading of its own, running
% straight from the abstract. Roughly 600–800 words, four or five
% paragraphs. Accessible to a reader from either side: a machine
% learning reader who has never heard of the secretory phase, and a
% clinician who has never trained an autoencoder.
%
% HIGH-LEVEL ONLY. Paragraph-by-paragraph breakdown comes after the
% flow below is approved.
% =====================================================================

\section*{Introduction}

\note{PARAGRAPH 1 — the disease, and why it is a computational problem
and not only a clinical one. Prevalence, diagnostic delay, the fact that
lesions are read by eye from biopsies. What a reader needs to know about
endometriosis to follow the rest, and nothing more. Participants are
people with a uterus; this is a human health problem.}

\note{PARAGRAPH 2 — the data situation. Eleven public studies, no shared
panel, no shared protocol, cohorts of tens of donors each. Why pooling
is the only route to statistical power here, and why pooling is also
what breaks the standard tooling.}

\note{PARAGRAPH 3 — THE CONFOUND, which is the intellectual core of the
paper and the thing a reader must understand before any result means
anything. Each donor was biopsied once. Menstrual phase is therefore
constant within a donor, and a per-donor offset is geometrically the same
object as a batch effect. Batch correction cannot tell them apart, so
removing the batch removes the cycle. State plainly that this is not a
technical caveat: it is the reason naive integration of endometrial data
silently discards the biology people want.}

\note{PARAGRAPH 4 — what we do about it. Not ``we propose a novel
architecture''. We make the trade-off explicit and measurable, map the
frontier between donor-blindness and retained cycle signal, and report
from a named point on it. Sixteen dimensions and 1.88M parameters are
mentioned here as a consequence of the design, not as a selling point.}

\note{PARAGRAPH 5 — what the paper contributes, as a short list the
Results then delivers in order. Flow rule: the last sentence here must
be the first question Results answers.}

\todo{deferred until \texttt{PUBLICATION\_PLAN} §0 reframing lands}
